% \begin{savequote}
% \includegraphics[width=0.7\linewidth]{./images/noauth/AASdev20090228-img110.jpg}
% \end{savequote}
\def\ChapterCode{KHD}
\chapter
[\CO{[KHD]} Einführung in den Gro{\ss}schadens- und Katastrophenhilfsdienst (.at)]
{Einführung in den Gro{\ss}schadens- und Katastrophenhilfsdienst (.at)}
\label{chp:KHD}

\ChapterIntro
{%
Großschadens- und Katastrophenfälle stellen die mit der
notfallmedizinischen Versorgung beauftragten Organisationen
vor nicht alltägliche Herausforderungen. 
Massenanfälle von Verletzten oder Erkrankten, zu wenig
Resourcen oder zerstörte Infrastruktur erfordern gänzlich
andere Herangehensweisen als im regulären Betrieb. 

}
{%
\begin{description}
  \item[Maintainer:] Roman Koch
  \item[Autoren:] Diverse
  \item[Reviewer:] Gerald Höritzmiller
  \item[Version:] {\VarDocVersionTypeName} ({\VarDocVersionTypeDescription})
  \item[SHA1:] {\DocGitCommit}
\end{description}

{\TxTeilkapitelAASS}
}

\section{Katastrophen, Großschadensereignisse, Unfall}



\begin{PictureSeriesWide}{}
{Bilderserie: Großschadensereignisse}
\PictureSeriesRowWide{3}
% {./images/noauth/AASdev20090228-img110.jpg}
% {ICE-Unglück von Eschede am 3. Juni
% 1998: 101 Tote, 
% 88 Schwerverletzte, 
% 106 Leicht- und Unverletzte.\SourcePicture{WMC (Nils
% Fretwurst)/PD}}
{./images/gabriel-sebastian-aass/UE2011FORTUNA-00942-0154pt-0300dpi}
{Bewältigung eines Großschadensfalles: Ordnung in das Chaos
bringen.}
{./images/gabriel-sebastian-aass/UE2011FORTUNA-00568-0154pt-0300dpi}
{Führungsrolle bie der Bewältigung eines
Großschadensereignis: Der Leiter einer Sanitätshilfestelle (SanHiSt). \SourcePicture{Gabriel}} {./images/gabriel-sebastian-aass/UEKAISERMUEHLENTUNNEL2011-000253-0154pt-0300dpi}
{Zusammenarbeit verschiedener Organisationen zur Bewältigung einer Schadenslage. \SourcePicture{Gabriel}}
{empty}
{}
\end{PictureSeriesWide}





\subsection{Grundlegende Begriffe}

\TextSlide{Unfall (Bagatelle, Notfall)}{%
Ein Unfall ist plötzlich und unerwartet, die Dauer des
Einsatzes beträgt in etwa Minuten bis zu einer Stunde. Die
\E{Infrastruktur}\footnote{\emph{Infrastruktur:} Wichtige
Basiseinrichtungen und Dienstleistungen, die das öffentliche
Leben ermöglichen; z.\,B. Straßen, Elektrizität, Telefon,
Krankenhäuser, Verwaltungseinrichtungen, Nahrungs- und
Wasserversorgung, Abfall- und Abwasserentsorgung, usw.} ist
\E{intakt}. Es sind bis zu ca. 14
Personen betroffen (Richtwert).

Die Erstversorgung erfolgt innerhalb von Minuten durch den regulären
Rettungsdienst nach den Regeln der \E{Individualmedizin}.
Der Abtransport wird im Routinebetrieb mit örtlich
vorhandenen Mitteln durchgeführt. Die definitive Versorgung
erfolgt in umliegenden Krankenhäusern, ebenfalls nach den
Regeln der \E{Individualmedizin}. Das aufgewendete Material
und das Personal stammt aus dem Normalbetrieb des Rettungs-
und Krankentransportes.
}
{\FullPrettyRef{tab:uebersicht-unfall-grossunfall-katastrophe}}
{}{}{}




\TextSlide{Großunfall (Großschadensereignis)}{
\index{Großunfall}
\index{Großschadensereignis}
\label{topic:grossschaden}
Ein Großunfall liegt vor, wenn
anzunehmen ist, dass das Ereignis \E{mit den örtlichen
personellen und materiellen Kräften und Mitteln nicht
bewältigt werden kann}, aber \E{keine erklärte
Katastrophensituation} vorliegt.

Das Ereignis ist \E{lokal begrenzt} und dauert mehrere
Stunden, die \E{Infrastruktur} ist jedoch \E{intakt}. Die
Erstversorgung von ungefähr 15 oder mehr Patienten
(Richtwert) erfolgt innerhalb von Minuten bis Stunden durch
den Rettungsdienst nach den Regeln der
\E{Katastrophenmedizin}. Der Abtransport erfolgt relativ
unbehindert nach der Erstversorgung, es sind genügend
Transportmittel vorhanden.

Die definitive Versorgung erfolgt in umliegende Krankenhäuser
nach den Regeln der \E{Individualmedizin}. Das Personal setzt
sich aus allen verfügbaren Helfern zusammen, welche sowohl aus
dem Normalbetrieb des Rettungs- und Krankentransportes und
von Sondereinsatzgruppen stammt. Es wird zusätzliches
Material eingesetzt.
}{\FullPrettyRef{tab:uebersicht-unfall-grossunfall-katastrophe}}{}{}{}




\TextSlide{Katastrophe (außergewöhnliche Lage)}{
\index{Katastrophe}\label{topic:katastrophe}
Eine Katastrophe\ZFootnote{Katastrophe vom \Lang{griech.}: 
schweres Unglück, Zusammenbruch,
Wendung zum Schlimmen \cite{BertelsmannVolkslexikon25,BrockhausDrei2Bd2}} 
ist ein Ereignis mit einem
\EE{außergewöhnlichen Schadensausmaß}, welches eine \EE{konkrete
Gefahr oder Schaden für Menschen}\Z{, Tiere, Umwelt, Kulturgüter} und
Sachwerte sowie für die \EE{Infrastruktur} zur Sicherstellung
der Versorgung mit lebensnotwendigen Gütern und
Dienstleistungen mit sich bringt. Es bedarf der \EE{koordinierten Führung durch
die Behörde}.\ZFootnote{Die genaue Definition der
Katastrophe unterscheidet sich je nach den jeweiligen
Katastrophenhilfegesetzen der Bundesländer.}

Das Ereignis betrifft oft ein \E{(über-) regionales Gebiet} und
kann \E{Tage bis Wochen} andauern. Die \E{Infrastruktur} ist
\E{gestört} oder sogar zerstört. Die Anzahl der Verletzten
variiert je nach Art des Ereignisses. Die Erstversorgung
erfolgt innerhalb von Stunden bis Tagen durch die
KHD\footnote{KHD: Katastrophen- und
Großschadenshilfsdienst}-Einheiten nach den Regeln der
\E{Katastrophenmedizin}. Der Abtransport ist stark behindert
und verzögert, es gibt nur wenige oder keine vorhandene
Transportmittel.

Die definitive Versorgung erfolgt vorwiegend auf
Verbandsplätzen, in Behelfsspitälern, oder in Lazaretten nach
den Regeln der \E{Katastrophenmedizin}. Das Material und 
Personal setzt sich zusammen aus Überlebenden und verfügbaren
Helfer, einsatzbereiten Fahrzeugen, Kat-Einsatz aller
Hilfsorganisationen, überregionaler und internationaler
Hilfe.

\subparagraph{Arten von Katastrophen}
\begin{itemize}
  \item \T{Naturkatastrophen} (\I{Elementarkatastrophen}):
  Erdbeben, Überschwemmungen, Erdrutsche, Lawinen, etc.
  \item \T{Technische Katastrophen}: Flugzeugabsturz,
  Eisenbahnunglück, Re\-aktor\-un\-fall, Schiffs\-untergang
  \item \T{Sekundärkatastrophen}: Hungersnot, Seuchen etc.
\end{itemize}
}{\FullPrettyRef{tab:uebersicht-unfall-grossunfall-katastrophe}}{}{}{}


% \TableWide
% 	{<Titel>}%1
% 	{<Beschreibung>}%2
% 	{<Label>}%3
% 	{<Quelle>}%4
% 	{<Lizenz>}%5
% 	{<Tabelle>}%6
% 	{<Float>}%8
\TableWide
{Übersicht Unfall\,--\,Grossunfall\,--\,Katastrophe}%1
{}%2
{}%3
{}%4
{}%5
{%
\label{tab:uebersicht-unfall-grossunfall-katastrophe}%
\small
\begin{tabularx}{\linewidth}{p{8em}XXX}
\toprule\addlinespace
\noindent
&
\noindent\TabHeading{Unfall}
&
\noindent\TabHeading{Großunfall}
&
\noindent\TabHeading{Katastrophe}
\\\addlinespace\midrule\addlinespace
\noindent\TabSub{Ereignis:}
 &
plötzlich, unerwartet
&
lokal begrenzt
&
(über)-regionales Gebiet
\\\addlinespace
\noindent\TabSub{Dauer:}
 &
Minuten bis 1 Std.
&
mehrere Stunden
&
Tage, Wochen
\\\addlinespace
\noindent\TabSub{Infrastruktur:}

&
intakt 
&
intakt
&
gestört, zerstört
\\\addlinespace
\noindent\TabSub{Patienten:}
 &
{\textless} $\sim$\,15 (Richtwert)
&
$\geq$ $\sim$\,15 (Richtwert)
&
keine bis viele, je nach Art des Ereignisses
\\\addlinespace
\noindent\TabSub{Erstversorgung:}
 &
innerhalb von Minuten durch Rettungsdienst nach Regeln der
Individual\-medizin
&
innerhalb von Minuten bis Stunden durch den Rettungsdienst nach den
Regeln der Katastrophenmedizin
&
innerhalb von Stunden bis Tagen durch KHD\footnotemark{}{}-Einheiten
nach den Regeln der Katastrophenmedizin
\\\addlinespace
\noindent\TabSub{Abtransport:}
 &
unbehindert, Routinebetrieb mit örtlich vorhandenen Mitteln
&
relativ unbehindert nach der Erstversorgung, genügend vorhandene
Transportmittel
&
stark behindert verzögert, wenige/nicht vorhandene Transportmittel
\\\addlinespace
\noindent\TabSub{Definitive Versorgung:}
 &
in umliegende Krankenhäuser nach Regeln der Individualmedizin
&
in umliegende Krankenhäuser nach den Regeln der Individualmedizin
&
auf Verbandsplätzen in Behelfsspitälern, in Lazaretten nach
Regeln der Katastrophenmedizin
\\\addlinespace
\noindent\TabSub{Material und Mann\-schaft:}
 &
Normalbetrieb Rettungs- und Krankentransport
&
alle verfügbaren Helfer, Normalbetrieb Rettungs- und
Krankentransport inkl. Sondereinsatzgruppe und zusätzlichem
Material
&
\"Uberlebende und verfügbare Helfer, einsatzbereite Fahrzeuge,
Kat-Einsatz aller Hilfsorganisationen, überregionale und
internationale Hilfe
\\\bottomrule
\end{tabularx}
\footnotetext{KHD: Katastrophenhilfsdienst}
}%6
{H}%8




\TextSlide{Individualmedizin}{%
Unter \T{Individualmedizin} versteht man die bedarfsgerechte, 
maximale und bestmögliche Versorgung jedes
Patienten.
Es besteht freie Arztwahl, welche auf das Vertrauensverhältnis zwischen
Patient und Arzt ausgerichtet ist. \bigskip
}{%
\begin{itemize}
  \item Bedarfsgerechte, maximal- und bestmögliche Versorgung jedes
Patienten, Ressourcen werden für den Einzelnen ausgeschöpft
\end{itemize}
}{}{}{}



\TextSlide{Katastrophenmedizin}{%
Unter \T{Katastrophenmedizin} versteht man die 
\EE{Massenversorgung} von Verletzten und Erkrankten
mit \EE{beschränkten Mitteln} und dem Zwang zur \EE{Einteilung nach
Dringlichkeit} und \EE{Möglichkeit} zur Behandlung.

Eine Vielzahl von Hilfsbedürftigen stehen wenigen Ärzten,
Sanitätern und sonstigen Helfern mit oft nicht ausreichenden
Mitteln gegenüber. Ziel muss es sein, unter den gegebenen
Umständen so viele Leben wie möglich zu retten, \E{mit
bewusstem Verzicht auf zeit- und personalaufwendige
Maximalversorgung des Einzelnen zugunsten der
lebensrettenden Minimalversorgung Vieler.}

Ein wesentliches Element der Katastrophenmedizin ist die
Einteilung der Patienten gemäß der Dringlichkeit und
Durchführbarkeit ihrer Behandlung oder ihres Transportes.
Diese Einteilung nennt man Triage (\FullPrettyRef{topic:triage}).
}{%
\begin{itemize}
  \item Aus beschränkten Mitteln das Beste für Viele
  herausholen, wobei auf maximale Versorgung Einzelner
  bewusst verzichtet werden muss.
\end{itemize}
}{}{}{}























\bigskip

% \clearpage
\section{Grundsätze zur Bewältigung von Großschadenslagen}



\TextSlide{\TxBeschreibung}{
Ziel ist es, das \EE{Bestmögliche} für die größte Anzahl an Patienten zur
rechten Zeit am rechten Ort zu erreichen. Das ist nur
erreichbar, wenn  der \E{Bestgeeignete} auf jeder Stufe
führt, klare und eindeutige \EE{Führungsstrukturen}
vorhanden sind, alle Einsatzkräfte ihre \E{Aufgabe kennen}, 
die \EE{Führungspyramide} von allen strikt eingehalten wird,
das dazu benötigte \E{Material vorhanden}, bereitgestellt
und rasch einsatzbereit ist, die Organisation eingespielt
und eingeübt ist, sowie Einsätze -- soweit möglich -- geplant
und vorbereitet sind .

Dafür wird zusätzliche Sonderausrüstung erforderlich, da ein
erhöhter Bedarf an Material, Medikamenten und Personal
besteht.  Man muss mit einem großen Patientenanfall, 
besonderen Verletzungsmustern, einer längeren
Einsatzzeit vor Ort und mit einer
Überforderung der Helfer rechnen.
}{
\TabSub{Ziele}
\begin{itemize}
    \item Das Bestmögliche, für die größte Anzahl,  zur rechten Zeit, am
    rechten Ort.
\end{itemize}

\TabSub{Voraussetzungen}
\begin{itemize}
    \item Führung durch Bestgeeigneten
    \item Führungsstrukturen, -pyramide
    \item Kenntnis der Aufgaben
    \item Material vorhanden und EB
    \item Übung, Plan und Vorbereitung
\end{itemize}

\TabSub{Sonderausrüstung} da
erhöhter Bedarf an:
\begin{itemize}
    \item Material
    \item Medikamenten
    \item Personal
\end{itemize}

\TabSub{Wir müssen rechnen mit}
\begin{itemize}
    \item großem Patientenanfall
    \item besondere Verletzungsmuster
    \item längere Einsatzzeit
    \item Überforderung der Helfer
\end{itemize}
}{}{}{}


\Text{Grundsätze des Einsatzes}
{%
\label{topic:einsatzgrundsaetze}%
Ziel ist es, ohne überstürztes Handeln die Ausdehnung der
Krisenlage auf benachbarte und andere Räume zu verhindern.
Dabei wird gemäß der gesetzlichen Bestimmungen von Kräften
der Einsatzorganisationen, Behörden und anderen
Einrichtungen koordiniert vorgegangen.
}

\begin{Synopsis}
  \SynopsisItem{Grundkomponenten eines KHD-Einsatzes: Führung -- Zeit -- Raum}
\end{Synopsis}

\Text{Führung}
{%
\label{topic:katastrophe-fuehrung}%
\DefinitionInPara{Als Führung bezeichnet man das steuernde
Einwirken auf andere Menschen um ein bestimmtes Ziel zu erreichen. }
Die Führung ist ein unentbehrliches Steuerungsinstrument zur
Erfüllung des Einsatzauftrages
(\enquote{Ziel}) und dem richtigen
Einsatz von Personal, Material und Transportmitteln. Im Schadensraum gilt:
\enquote{\EE{Ein einziger Verantwortlicher}} mit
Entscheidungsbefugnis und den nötigen Führungsmittel
\E{(\enquote{Einheit der Führung})}.
Dieser wirkt mittels möglichst einfacher Ma{\ss}nahmen
\E{(\enquote{Einfachheit})} auf den
Einsatzverlauf ein
\E{(\enquote{Handlungsfreiheit})}.
Mittel und Mannschaften werden dabei gemäß ihrer Eigenart
und Leistungsfähigkeit eingesetzt
\E{(\enquote{Ökonomie der Kräfte})}.
Er muss rasch auf Änderungen der Situation reagieren
\E{(\enquote{Beweglichkeit})} und
gleichzeitig vorausschauend handeln
\E{(\enquote{Reservenbildung})}.
}

\begin{Danger}
  \item \enquote{Führen} und
  \enquote{Behandeln} ist nicht gleichzeitig
  möglich!
\end{Danger}

\Text{Sanitäter und Notärzte müssen Führungsaufgaben übernehmen}
{%
Der erste am Schadensort eintreffende Sanitäter muss:

\begin{enumerate}
    \item Lagemeldung abgeben
    \item Erste, vorläufige Grundstrukturen einer SanHiSt
    aufbauen (provisorischer Einsatzleiter,
    festlegen von Sammelstellen, Wagenhalteplatz, etc.)
    \item Mit bereits anwesenden Einsatzkräften die ersten
    Maßnahmen zur Rettung und Behandlung festlegen
\end{enumerate}
}

\Text{Zeit}
{%
\label{topic:katastrophe-phasen}%
Jedes Großschadens- oder Katastrophenereignis zeigt unabhängig von seiner
Ursache einen charakteristischen Zeitablauf, vgl.
\FullPrettyRef{tab:khd-einsatz-phasen}.}

\begin{table}[H]\FontTableBase
\Captionof{table}{Phasen des Einsatzes.
\label{tab:khd-einsatz-phasen}}

\begin{tabularx}{\linewidth}{XXX}\toprule\addlinespace
\noindent\TabHeading{1. Phase}
 &
\noindent\TabHeading{2. Phase}
 &
\noindent\TabHeading{3. Phase} 
\\\midrule\addlinespace
Isolation (Minuten)
 &
Retten (Stunden)
 &
Wiederherstellung (Tage)
\\\addlinespace
Schaden begrenzen  

- Lage erkennen 

- Melden
 &
Ordnung schaffen 

- Kommandoposten
 &
Lokale Einsatzleitung 

- Verbindung 

- Lage beurteilen
\\\addlinespace
Überleben
 &
Weiterleben
 &
Endbehandlung
\\\addlinespace
Laien: 

- Nothilfe/Selbsthilfe 

- Retten / Erste Hilfe
 &
Geschulte Helfer: 

- Lebensrettende Sofortmaßnahmen 

- Transportfähigkeit
 &
Spezialisten: 
- Spital nach Dringlichkeit behandeln
\\\addlinespace
\multicolumn{3}{c}{\includegraphics[width=0.5\linewidth,height=1em]{./images/khd/arrow-right.png}}
\\\addlinespace\bottomrule
\end{tabularx}
\end{table}



\section{Katastrophenschutz}


\subsection{Katastrophenhilfegesetze regeln die Zuständigkeiten}
\label{topic:katastrophenhilfsgesetze}

Die Organisation der Katastrophenhilfe fällt in den Aufgabenbereich der Bundesländer bzw. der Gemeinden.
Dementsprechend verfügt jedes Bundesland über ein inhaltlich abweichendes 
Rettungsdienst- und Katastrophenschutzgesetz, welches die Aufgaben der Behörden und
Organisationen, sowie Zuständigkeiten in deren Bereich regelt.
\EE{Ausschließlich eine Behörde erklärt ein
außergewöhnliches Schadensereignis zur Katastrophe}. 
Behördliche Alarmpläne dienen zur Festlegung
der Maßnahmen der Katastrophenhilfe und zur Koordination der
im Anlassfall tätig werdenden Einrichtungen. Je nach
Ausdehnung des Schadensgebietes gibt es unterschiedliche
Kompetenzen:

\begin{center}
\TableBaseModification
\begin{tabulary}{\linewidth}{lclp{5cm}}
\toprule
\E{Gemeinde} 
 &
 ${\rightarrow}$
 &
Bürgermeister
&
~\TakeNotesMedium~
\\
\E{Bezirk} 
 &
 ${\rightarrow}$
 &
Bezirkshauptmann
&
~\TakeNotesMedium~
\\
\E{Bundesland}
 &
 ${\rightarrow}$
 &
Landeshauptmann
&
~\TakeNotesMedium~
\\
\E{Bundesgebiet}
 &
 ${\rightarrow}$
 &
Innenminister%, Bundeskanzler 
&
~\TakeNotesMedium~
\\\bottomrule
\end{tabulary}
\end{center}


\subsection{Die Katastrophenwarnung mittels Sirene alarmiert die
Bevölkerung}
\label{topic:alarmierung}

Österreich verfügt über ein flächendeckendes,
betriebsbereites Netz von 8\,120 Sirenen (Stand 10/2009).
Einmal wöchentlich erfolgt um 12:00 Uhr eine Sirenenprobe
mittels eines 15-sekündigen Dauertons.
Weiters  findet einmal jährlich am ersten Samstag im
Oktober zwischen 12:00 und 13:00 Uhr eine österreichweite
Sirenenprobe mit allen Katastrophensignalen statt.
\cite{wiki:zivilschutzsignale20120810}

\FullPrettyRef{fig:zvivilschutz-sirenensignale-at}, gibt eine
Übersicht über die in Österreich verwendeten Sirenensignale

\begin{PictureSeriesWide}{}{Zivilschutz-Sirenensignale in
Österreich\label{fig:zvivilschutz-sirenensignale-at}}
\PictureSeriesRowWide{4}
{./images/wikipedia-sh/AT-Sirene-Warnung}
{\EE{Warnung}: Herannahende Gefahr wird mit einem
dreiminütigen Dauerton angekündigt. Die Bevölkerung wird
damit aufgefordert Radio- oder Fernsehgerät einzuschalten
und dort bekanntgegebene Anordnungen zu beachten.
% \cite{wiki:zivilschutzsignale20120810}
}
{./images/wikipedia-sh/AT-Sirene-Alarm}
{\EE{Alarm}: Das Alarmsignal besteht aus einem auf- und
abschwellenden Ton von einer Minute Dauer und bedeutet unmittelbare Gefahr:
schützende Räumlichkeiten aufsuchen, über Medien durchgegebene
Verhaltensmaßnahmen befolgen und Radio
einschalten. 
% \cite{wiki:zivilschutzsignale20120810}
}
{./images/wikipedia-sh/AT-Sirene-Entwarnung}
{\EE{Entwarnung}: Das Ende der Gefahr wird mit einem
einminütigen Dauerton angezeigt. Mögliche Einschränkungen im täglichen Lebenslauf
werden über Medien
durchgegeben.
% \cite{wiki:zivilschutzsignale20120810}
}
{./images/wikipedia-sh/AT-Sirene-Sirenenprobe}
{\EE{Sirenenprobe}: Die Sirenenprobe ist ein 15 Sekunden
dauernder Dauerton oder ein besonders kurzer Alarmton, wobei die Sirene nur einen
leisen, kurzen Heuler von sich gibt, oder der normale
Feuerwehralarmton.
% \cite{wiki:zivilschutzsignale20120810}
}
% 
% \PictureSeriesRow{3}
% {./images/wikipedia-sh/AT-Sirene-Feuerwehralarm}
% {\EE{Feuerwehralarm}: Der Feuerwehralarm besteht aus dreimal
% 15 Sekunden Dauerton mit zweimal 7 Sekunden Unterbrechung. Die Verwendung des
% Alarmsignals hängt von der einzelnen Feuerwehr ab.
% \cite{wiki:zivilschutzsignale20120810}} 
% {./images/wikipedia-sh/AT-Sirene-Sirenenprobe}
% {\EE{Sirenenprobe}: Die Sirenenprobe ist ein 15 Sekunden
% dauernder Dauerton oder ein besonders kurzer Alarmton, wobei die Sirene nur einen
% leisen, kurzen Heuler von sich gibt, oder der normale
% Feuerwehralarmton. \cite{wiki:zivilschutzsignale20120810}}
% {empty}
% {}
% {empty}
% {}

\end{PictureSeriesWide}

% \begin{WidePar}
% \includegraphics[width=0.24\linewidth]{./images/noauth/khd-sirenen.png}\hfill\includegraphics[width=0.24\linewidth]{./images/noauth/khd-sirenen.png}\hfill\includegraphics[width=0.24\linewidth]{./images/noauth/khd-sirenen.png}\hfill\includegraphics[width=0.24\linewidth]{./images/noauth/khd-sirenen.png}
% 
% 
% \begin{figure}[H]
% \dotfill
% 
% \begin{WidePar}
% 
% \includegraphics[width=0.24\linewidth]{./images/noauth/khd-sirenen.png}\hfill\includegraphics[width=0.24\linewidth]{./images/noauth/khd-sirenen.png}\hfill\includegraphics[width=0.24\linewidth]{./images/noauth/khd-sirenen.png}\hfill\includegraphics[width=0.24\linewidth]{./images/noauth/khd-sirenen.png}
% 
% \subfloat
% [text]
% {\includegraphics[width=0.24\linewidth]{./images/noauth/khd-sirenen.png}}
% \subfloat
% [text]
% {\includegraphics[width=0.24\linewidth]{./images/noauth/khd-sirenen.png}}
% \subfloat
% [text]
% {\includegraphics[width=0.24\linewidth]{./images/noauth/khd-sirenen.png}}
% \subfloat
% [text]
% {\includegraphics[width=0.24\linewidth]{./images/noauth/khd-sirenen.png}}
% 
% \includegraphics[width=0.24\linewidth]{./images/noauth/khd-sirenen.png}\hfill\includegraphics[width=0.24\linewidth]{./images/noauth/khd-sirenen.png}\hfill\includegraphics[width=0.24\linewidth]{./images/noauth/khd-sirenen.png}\hfill\includegraphics[width=0.24\linewidth]{./images/noauth/khd-sirenen.png}
% \end{WidePar}
% \end{figure}
% \includegraphics[width=0.24\linewidth]{./images/noauth/khd-sirenen.png}\hfill\includegraphics[width=0.24\linewidth]{./images/noauth/khd-sirenen.png}\hfill\includegraphics[width=0.24\linewidth]{./images/noauth/khd-sirenen.png}\hfill\includegraphics[width=0.24\linewidth]{./images/noauth/khd-sirenen.png}
% 
% \end{WidePar}


\cite{wiki:zivilschutzsignale20120810}

% \subsection{Katastrophenschutz in Wien}
% \label{topic:katastrophenschutz}
% \TODO{EMHO}{2010-08-04}{Katastrophenhilfseinheiten fehlen!}{}
% \begin{center}
% \begin{tabular}{p{4.3190002cm}m{8.981cm}}
% Wien ist für den Groß\-schadens- und Katastrophenfall gut
% gerüstet. Tausende Helfer sind dazu im
% Katastrophen\-schutz-Kreis
% (\emph{\enquote{K-Kreis}})
% org\-anisiert.
% 
% Der K-Kreis ist ein Zu\-sammen\-schluss der beruflichen und
% freiwilligen Hilfs- und Ein\-satz\-organisationen sowie
% wichtiger Behörden. Dreh\-scheibe und gemeinsame
% Anlaufstelle ist dabei die Stadt Wien, Magistrats\-direktion
% - Katastrophen\-manage\-ment und Sofort\-ma{\ss}\-nahmen und
% die Helfer Wiens.
%  &
% \includegraphics[width=9.046cm,height=13.623cm]{./images/noauth/k-kreis08.png}
% \CaptionofDiff{figure}{Der K-Kreis.
% \SourcePicture{http://www.diehelferwiens.at/k-kreis , Copyright: Die Helfer
% Wiens, Grafik: Erdle/Zimmermann/Schimanek}}{Der K-Kreis}
% \\
% \end{tabular}
% \end{center}

\subsection{Katastrophenhilfe}

Hilfsmaßnahmen, die von einer Sanitätsorganisation erwartet
werden, umfassen:

\begin{itemize}
\item Bergen und Retten, Versorgung und Transport 
\item Pflegerische Betreuung 
\item Sorge für Unterkunft und Verpflegung 
\item Karitative Betreuung von Hilfsbedürftigen
\end{itemize}

\begin{center}
\FontTableBase
\begin{tabularx}{\linewidth}{lX}
\toprule
\TabHeading{Sanitätseinsatz}
 &
\begin{itemize}
     \item Retten/Bergen
     \item Triage
     \item Lebensrettende Sofortmaßnahmen
     \item Erstversorgung
     \item Sanitätshilfe
     \item Abtransport
\end{itemize}
\\\midrule
\TabHeading{Pflegeeinsatz}
&
 Wenn die pflegerische Versorgung durch die bestehenden
 Einrichtungen nicht mehr sichergestellt ist.
 \begin{itemize}
     \item Erweiterung der vorhandenen Behandlungsmöglichkeiten
     \item Errichtung von Hilfskrankenhäusern
 \end{itemize}
\\\midrule
\TabHeading{Sozialeinsatz}
&
Wenn Menschen infolge einer Katastrophe ihre
Selbstversorgungsmöglichkeit verloren haben.

\begin{itemize}
    \item Einrichtung/Betrieb von Notunterkünften/Lagern
    \item Beschaffung, Zubereitung und Ausgabe von Lebensmitteln, Kleidern,
    Gebrauchsgegenständen
    \item Fürsorgliche Hilfe, Registrierung und Betreuung
\end{itemize}
\\\bottomrule
\end{tabularx}
\end{center}



\clearpage
\section{Triage und Triagegruppen}
\label{topic:triage}
\index{Triage}


\Text{\TxBeschreibung}
{%
\DefinitionInPara{Als \T{Triage} bezeichnet
man das \EE{Begutachten} aller Patienten mit geringen
Mitteln und Aufwand sowie das \EE{Einteilen} in Triagegruppen 
gemäß der \EE{Behandlungsdringlichkeit} und \EE{-möglichkeit} mit dem Ziel mit
den vorhandenen Resourcen soviel als möglich zu erreichen. \Z{Triage
\Lang{franz.}: Sichtung, Auswahl.}}
Anhand der Triage können Prioritäten bei der Behandlung von 
Patienten gesetzt werden.
Die Tabelle
\FullPrettyRef{tab:khd-behandlungsstellen-uebersicht}, gibt eine
Übersicht über die Triagegruppen.
}

\begin{Synopsis}
  \SynopsisItem{Grundsatz: 
  \Speech{\enquote{Life before Limbs}}: 
  Leben vor Gliedmaßen, 
  Funktion vor Schönheit!}
\end{Synopsis}

\TextSlide{Triagegruppen können sich im Verlauf ändern}{%
Die Zuordnung zu einer Triagegruppe kann sich im Verlauf
verändern z.\,B. wenn sich der Zustand des Patienten ändert
oder mehr Resourcen zur Verfügung stehen. Eine ständige
\E{Nachtriage} ist daher erforderlich! 
Typische Beispiele für nachträgliche Änderungen der Triagegruppe wären:
\begin{enumerate}
    \item Verschlechterung des Patientenzustandes: 
    \texttt{III} $\rightarrow$ \texttt{I}
    \item Verbesserung des Patientenzustandes nach erfolgter Behandlung: 
    \texttt{I} $\rightarrow$ \texttt{III}
    \item Verbesserung der Versorgungslage 
    (zusätzliches Personal bzw. Material verfügbar, 
    oder Behandlungskapazitäten sind frei geworden): 
    \texttt{IV} $\rightarrow$ \texttt{I}
\end{enumerate}
}{

}{}{}{}


\TextSlide
{Transportpriorität}
{%
\label{topic:transportprioritaet}
Nach einer erfolgten Behandlung kann eine \T{Transportpriorität} 
zugewiesen werden. 
Man unterscheidet hierbei: 

\begin{itemize}
  \item \EE{Transportpriorität A}: Hohe Transportpriorität. 
  Der rasche Transport zur
  frühzeitigen Fachbehandlung nach ärztlicher Notversorgung ist
  angesagt.
  \item \EE{Transportpriorität B}: Niedrige Transportpriorität. 
  Der Transport zur Fachbehandlung
  kann verzögert erfolgen.
\end{itemize}
Die Transportpriorität ist \E{unabhängig von der Triagegruppe}! \Z{Als Richtwert kann die Notwendigkeit zur operativen Versorgung herangezogen werden: 
Verletzungsbilder, 
welche innerhalb von 6 Stunden operativ (bzw. definitiv) versorgt werden müssen, 
erhalten eher Transportpriorität A, 
alle anderen Verletzungsbilder Transportpriorität B.}
}
{
\begin{itemize}
  \item \EE{A}: Rascher Transport zur
  frühzeitigen Fachbehandlung
  \item \EE{B}: Transport zur Fachbehandlung
  kann verzögert erfolgen
  \item Unabhängig von Triagegruppe
\end{itemize}
}
{}
{}
{}


% \TableWide
% 	{<Titel>}%1
% 	{<Beschreibung>}%2
% 	{<Label>}%3
% 	{<Quelle>}%4
% 	{<Lizenz>}%5
% 	{<Tabelle>}%6
% 	{<Float>}%8
\TableWide
{Übersicht: Triagegruppen}%1
{}%2
{}%3
{}%4
{}%5
{%
\label{tab:khd-behandlungsstellen-uebersicht}\label{tab:khd-triagegruppen-uebersicht}
\noindent\begin{tabularx}{\linewidth}{XXXX}\addlinespace\toprule\addlinespace
\noindent\TabHeading{I: Sofortbehandlung vor Ort}
&
\noindent\TabHeading{II: Dringliche Behandlung}
&
\noindent\TabHeading{III: Warten Leichtverletzte}
&
\noindent\TabHeading{IV: Warten Schwerverletzte}
\\\addlinespace\midrule\addlinespace
Lebensrettende Soforteingriffe
&
Schwerverletzte
&
Leichtverletzte
&
Schwerverletzte
\\\addlinespace
\includegraphics[width=3cm]{./images/khd/KHD-Symbole-RK/pdf/BH1}
&
\includegraphics[width=3cm]{./images/khd/KHD-Symbole-RK/pdf/BH2}
&
\includegraphics[width=3cm]{./images/khd/KHD-Symbole-RK/pdf/BH3}
&
\includegraphics[width=3cm]{./images/khd/KHD-Symbole-RK/pdf/BH4}
\\\addlinespace\midrule\addlinespace
\multicolumn{4}{c}{\TabSub{Welche Patienten?}}
\\\addlinespace\midrule\addlinespace
Patienten, welche ohne sofortige Behandlung sterben bzw. schwere
Schäden erleiden
&
Behandlung und Herstellung der Transportfähigkeit von schwer
verletzten/erkrankten Personen, eine weitere Behandlung ist
nur mehr im Krankenhaus möglich.
&
Patienten, welche nur einer Minimalbehandlung bedürfen.
&
Patienten, welche mit dem zur Verfügung stehenden Mitteln
(Personal, Material) \EE{derzeit nicht} versorgt werden können. Abwartende
Behandlung.
% \\\addlinespace\midrule\addlinespace
% \multicolumn{4}{c}{\TabSub{Zum Beispiel}}
% \\\addlinespace\midrule\addlinespace
% \begin{itemize}
%     \item[--] Atemstörungen
%     \item[--] Spannungs\-pneu\-mo\-thorax
%     \item[--] schwere äußere Blutungen
%     \item[--] schwerer traumatischer Schock
% \end{itemize}
% &
% \begin{itemize}
%     \item[--] SHT
%     \item[--] Wirbelverletzungen mit Lähmungen
%     \item[--] Bauchtrauma, innere Blutungen
%     \item[--] schwere Thoraxtraumen
%     \item[--] große Gefäßverletzungen
%     \item[--] offene Knochenbrüche und  Gelenks\-verletzungen
%     \item[--] Augenverletzungen
%     \item[--] große Weichteilverletzungen
%     \item[--] Verbrennungen \VonBis{15}{30}\PC*
%     \item[--] geschlossene Knochenbrüche und Gelenks\-ver\-letzungen
% \end{itemize}
% &
% \begin{itemize}
%     \item[--] Verbrennungen bis 15\PC*
%     \item[--] kleine Weichteilwunden
%     \item[--] einfache Knochenbrüche
%     \item[--] Prellungen, Zerrungen
%     \item[--] leichtes SHT
% \end{itemize}
% &
% \begin{itemize}
%     \item[--] schweres Polytrauma
%     \item[--] schwerstes \Abk{SHT}
%     \item[--] offene Körperhöhlentraumen
%     \item[--] Mehrhöhlenverletzungen
%     \item[--] Verbrennungen über 50\PC*
% \end{itemize}
\\\addlinespace\midrule\addlinespace
\multicolumn{4}{c}{\TabSub{Maßnahmen, Beispiele}}
\\\addlinespace\midrule\addlinespace
\begin{itemize}
    \item[--] Überprüfung der Triage
    \item[--] Volumenersatz
    \item[--] Intubation
    \item[--] Thoraxdrainage
    \item[--] Noteingriffe
\end{itemize}
&
\begin{itemize}
    \item[--] Überprüfung der Triage
    \item[--] Verbände
    \item[--] Schockbehandlung
    \item[--] Fixierung/Schienung
    \item[--] Schmerzbekämpfung
    \item[--] pflegerische Maßnahmen
    \item[--] Lebensrettende Sofortmaßnahmen
\end{itemize}
&
\begin{itemize}
    \item[--] Überprüfung der Triage
    \item[--] pflegerische Maßnahmen
    \item[--] Betreuung (Kriseninterventionsteam (\Abk{KIT})
    \item[--] Entlassung ambulant Behandelter zur \enquote{Sammelstelle Unverletzte}
\end{itemize}
&
\begin{itemize}
    \item[--] Überprüfung der Triage
    \item[--] Behandlung  (Schmerzbekämpfung)
    \item[--] Psychologische Betreuung, evtl. durch \Abk{KIT}
    \item[--] Überwachung
\end{itemize}
\\\addlinespace\bottomrule

\end{tabularx}
}%6
{H}%8

\clearpage
\section{Räume und Sanitätshilfsstelle (SanHiSt)}
\label{topic:sanhist}

\TextSlide{Übersicht}{
Der Einsatz wird unterteilt in: 
\begin{itemize}
  \item \T{Schadensraum}: Der Schadensraum ist der Ort des eigentlichen (Schadens-) Geschehens
  innerhalb der inneren Absperrung. 
  Es gibt in ihm eine einheitliche Organisations\-form für den
  Sanitätsdienst: 
  Die \T{Sanitätshilfsstelle} (\T{SanHiSt}), 
  welche aus den drei Räumen 
  \T{Triageraum}, 
  \T{Behandlungsraum} und
  \T{Transportraum} besteht.
  \item \T{Transportraum}: 
  Der Transportraum umfasst in diesem Zusammenhang alle Mittel 
  und Wege um
  die Patienten vom Schadens\-raum zum Hospitalisationsraum zu
  verbringen. 
  Dieser \enquote{allgemeine Transportraum} ist nicht zu ver\-wechseln mit dem 
  Transportraum der Sanitätshilfsstelle (s.\,u.).
  \item \T{Hospitalisationsraum}: 
  Der Hospitalisationsraum umfasst die Transportziele 
  und Aufnahmespitäler.
\end{itemize}
}{
\begin{itemize}
  \item Schadensraum $\rightarrow$ SanHiSt
  \begin{itemize}
    \item Triageraum
    \item Behandlungsraum
    \item Transportraum
  \end{itemize}
  \item Transportraum
  \item Hospitalisationsraum
\end{itemize}
}{}{}{}


 


% \TableWide
% 	{<Titel>}%1
% 	{<Beschreibung>}%2
% 	{<Label>}%3
% 	{<Quelle>}%4
% 	{<Lizenz>}%5
% 	{<Tabelle>}%6
% 	{<Float>}%8
\TableWide
{Gliederung des Einsatzraumes}%1
{}%2
{}%3
{}%4
{}%5
{%
\begin{Parallel}
{%
% \includegraphics[width=\linewidth]{./images/khd/sanhist-plan-edited.png}%
\includegraphics[width=\linewidth]{./images/koch-roman-aass/SANHIST}%
}
{%
\begin{itemize}
  \item Sicherheitseinrichtungen
  \begin{itemize}
    \item {\T[Absperrung!äußere@Absperrung!aeussere]{Äußere Absperrung}:} 
    Großräumige Verkehrsumleitung
    \item {\T[Absperrung!innere@Absperrung!innere]{Innere Absperrung}:} 
    hält Schaulustige fern, 
    hält Personen in Panik zurück
    \item \T{Sicherheitsring}: 
    Absperrung des Schadenplatzes, Sicherheit für Einsatzkräfte
    \item \T{Pforte}: Einbahnsystem für Einsatz-Kfz, keine Fremdfahrzeuge, 
    möglichst nur eine Pforte
  \end{itemize}
  \item Schadensraum
  \begin{itemize}
    \item \T{Sanitätshilfsstelle} (\T{SanHiSt}):
    \begin{itemize}
       \item Triageraum
       \item Behandlungsraum mit Behandlungsstellen \Range{1}{4}
       \item Transportraum mit Wagenhalteplatz, Verladestelle, ggfs. Hubschrauberlandeplatz
       \item Sammelstellen (außerhalb platziert):
       \begin{itemize}
         \item \T{Sammelstelle Tote} (t) und
         \item \T{Sammelstelle Unverletzte} (U)
       \end{itemize}
       \item \T{Material- und Meldestelle} (M)
       \item Betreuungsstelle: 
       Sie dient zur umfassenden psychosozialen Betreuung der Betroffenen (Krisenintervention (KIT)) und Mitarbeiter zuständig. 
    \end{itemize}
    \item \T{Bergetriage} (BT)
    \item Leitung und Infrastruktur:
    \begin{itemize}
      \item \T{Einsatzleiter} (EL)
      \item \T{Leitender Notarzt} (LNA)
      \item \T{Leiter Information} (i) 
      für Weiterleitung von Informationen an die Außenwelt (Presse, Angehörige, {\ldots})
    \end{itemize}
  \end{itemize}
  \item Hospitalisationsraum
\end{itemize}%

\url{http://www.roteskreuz.at/nocache/organisieren/organisation/wer-wir-sind/rechtliche-grundlagen/vorschriften/grossunfall/}

\cite{RkRahmenvorschriftGrossunfaelle2007}
\url{http://www.roteskreuz.at/nocache/organisieren/organisation/wer-wir-sind/rechtliche-grundlagen/vorschriften/grossunfall/do/download/uid/506/}
\cite{RkPlsNeu2008} \url{http://www.roteskreuz.at/nocache/organisieren/organisation/wer-wir-sind/rechtliche-grundlagen/vorschriften/grossunfall/do/download/uid/2815/}
}
\end{Parallel}
}%6
{H}%8



\Dictum
[Marie von Ebner-Eschenbach]
{Wer in die Öffentlichkeit tritt, hat keine Nachsicht zu erwarten und keine zu fordern.}

\clearpage
\subsection{Triageraum}



\TextSlideTeaser
{Bergetriage}%1 Titel
{%
Die Bergetriage nimmt die Festlegung der Bergepriorität vor,
d.\,h. in welcher Reihenfolge die Patienten zur Triagestelle
gebracht werden sollen. 
Die \T{Bergepriorität} kennt drei Gruppen:
\begin{enumerate} 
  \item Dringend
  \item Normal
  \item Verstorben
\end{enumerate} 
Gekenntzeichnet werden die Patienten mittels des
Patientenleitsystems (PLS).

Die Bergetriage wird nur gebildet wenn keine
Gefahrenzone und genügend Personal vorhanden ist.
Das Team muss mind. aus einem NFS und einem NA bestehen. 
Die Bergetriage behandelt nicht und betreibt keine
detaillierte Diagnostik!}%2 Text
{
\begin{itemize}
  \item Festlegung der Bergepriorität
  \begin{enumerate} 
    \item Dringend
    \item Normal
    \item Verstorben
  \end{enumerate} 
  \item Mind. NA + NFS
  \item Gefahren beachten
  \item Behandelt nicht, keine detaillierte Diagnostik
\end{itemize}}%3 Slide
{./images/khd/KHD-Symbole-RK/pdf/BT}%4 Bilddatei
{Symbol Bergetriage}%5 Bildtitel
{}%6 Bildbeschreibung
{---}%7 Quelle
{---}%8 Lizenz
 % Vormals \Layout{60}



\TextSlideTeaser
{Triage (Sichtung)}%1 Titel
{%
Aufgabe der Triagestelle ist die Zuordnung der Patienten zu den Triagegruppen
(\begin{inparadesc}
\item[\Code{I}] Sofortbehandlung (lebensrettende Soforteingriffe); 
\item[\Code{II}] Dringliche Behandlung (Schwerverletzte); 
\item[\Code{III}] Warten Leichtverletzte; 
\item[\Code{IV}] Warten Schwerverletzte
\end{inparadesc}) und den Sammelstellen. \E{Spätestens} hier erfolgt die
zentrale Registrierung aller Betroffenen mittels
Patientenleitsystem (PLS).

Die Triage erfolgt schnell, mit geringem Aufwand und unter
hohem Zeitdruck. Der maximale Zeitaufwand pro Patient
beträgt bei gehenden Patienten 1 Minute und bei liegenden
Patienten 3 Minuten. Es erfolgt \E{keine detaillierte}
Diagnostik und keine Behandlung.

% \begin{center}
% % \includegraphics[width=2.cm]{./images/noauth/sanhist-triage.png}
% \TODO{GABS}{2010-04-27}{Foto Triageplatz}{}
% \end{center}

Es werden zwei Triageplätze eingerichtet, damit sich der Notarzt nur zum
nächsten Patienten umdrehen braucht und abwechselnd triagieren kann. 
Die Vorbereitung des Patienten erfolgt durch Sanitätsfachpersonal und umfasst
das Entkleiden sowie die Vornahme des Ersteinschätzungsblockes. }%2 Text
{
\begin{itemize}
  \item Zentrale Registrierung aller Patienten mittels Patientenleitsystem
  (PLS)
  \item Zuordnung der Patienten zu Triagegruppen
  \VonBis{I}{IV} und Sammelstellen U und t
  \item Schnell, mit geringem Aufwand, hoher Zeitdruck.
  
  Maximaler Zeitaufwand pro Patient gesamt:
  \begin{itemize}
    \item Gehend 1 Minute
    \item Liegend 3 Minuten
  \end{itemize}
  \item Keine detaillierte Diagnostik
\end{itemize}}%3 Slide
{./images/gabriel-sebastian-aass/UE2011FORTUNA-00932-0154pt-0300dpi}%4 Bilddatei
{Leiter Triagestelle}%5 Bildtitel
{}%6 Bildbeschreibung
{Sebastian Gabriel}%7 Quelle
{MfG}%8 Lizenz
 % Vormals \Layout{60}










\subsection{Behandlungsraum}



\TextSlide{Behandlungsstellen}{%
Es werden vier Behandlungsstellen entsprechend der
Triagegruppen (\FullPrettyRef{tab:khd-behandlungsstellen-uebersicht}) eingerichtet, 
vgl. \FullPrettyRef{fig:behandlungsstellen}.
}{%
\begin{itemize}
  \item Analog zu Triagegruppen
  \item \FullPrettyRef{fig:behandlungsstellen}
\end{itemize}
}




\begin{PictureSeries}{}{Behandlungsraum und Behandlungsstellen: Die Kennzeichnung erfolgt analog zu den Triagegruppen.\label{fig:behandlungsstellen}}
\PictureSeriesRow{3}
{./images/khd/KHD-Symbole-RK/pdf/BH1}
{Sofortbehandlung vor Ort}
{./images/khd/KHD-Symbole-RK/pdf/BH2}
{Dringliche Behandlung}
{./images/khd/KHD-Symbole-RK/pdf/BH3}
{Warten Leichtverletzte}
{./images/khd/KHD-Symbole-RK/pdf/BH4}
{Warten Schwerverletzte}
\PictureSeriesRow{3}
{./images/khd/KHD-Symbole-RK/pdf/BH4}
{Warten Schwerverletzte}
{./images/khd/KHD-Symbole-RK/pdf/Behandlung}
{Kennzeichnung des Behandlungsraums bzw. des \enquote*{Leiter Behandlungsraum}}
{empty}
{}
{empty}
{}
\end{PictureSeries}

% \clearpage
\subsection{Transportraum}


\TextSlideTeaser
{Verladestelle}%1 Titel
{%
Hier erfolgt die Verladung von Patienten in die Fahrzeuge. 
Es muss dabei mittels einer Liste
protokolliert werden, welcher Patient mit welchem
Transportmittel in welches Krankenhaus eingeliefert wird. 
Außerdem  wird
der erste Abschnitt der PLS beim Verladen abgegeben.

Die Verladestelle sollte -- sofern möglich -- in der Nähe
der Behandlungsstelle II errichtet werden.
Es gibt eine definierte Ein- und Ausfahrt, eine
Einbahnregelung ist 
anzustreben! }%2 Text
{%
\begin{itemize}
  \item Protokollierung!
  \item Ein-, Ausfahrt, Einbahnsystem
\end{itemize}}%3 Slide
{./images/khd/KHD-Symbole-RK/pdf/Verladestelle}%4 Bilddatei
{Kennzeichnung Transportraum}%5 Bildtitel
{}%6 Bildbeschreibung
{---}%7 Quelle
{---}%8 Lizenz
 % Vormals \Layout{60}



\TextSlideTeaser
{Wagenhalte\-platz}%1 Titel
{%
Der Wagenhalteplatz muss \EE{frühzeitig} bestimmt werden, da
weitere Kräfte rasch nachfolgen werden. Er sollte in
ausreichender \EE{Entfernung} festgelegt werden um den
späteren Aufbau der weiteren SanHiSt nicht zu behindern.
Der Triage- und Behandlungsraum (und die
Verladestelle) benötigen viel Raum, um die erforderlichen
Fahrzeuge aufnehmen zu können. Wenn möglich gilt
Schrägparkordnung,  um ein freies Be- und Entladen, 
sowie Wegfahren zu ermöglichen.
Die Witterung und die Beschaffenheit des Untergrunds ist zu
beachten.
Die Fahrzeuge sollen nach Einsatzmittelkategorien
(\Abk{NAW}, \Abk{RTW}, \ldots) sortiert abgestellt werden.
Die Schlüssel müssen stecken gelassen werden!}%2 Text
{
\begin{itemize}
    \item Frühzeitig, Abstand von SanHiSt
    \item Groß genug; Bodenbeschaffenheit und Witterung bedenken!
    \item Schrägparkordnung, freies Heck zum Ent-/Beladen
    \item Sortiert nach Einsatzmittel
    \item Schlüssel stecken lassen
\end{itemize}}%3 Slide
{./images/khd/KHD-Symbole-RK/pdf/Transport_kfz}%4 Bilddatei
{Kennzeichnung Wagenhalte\-platz}%5 Bildtitel
{}%6 Bildbeschreibung
{}%7 Quelle
{}%8 Lizenz
 % Vormals \Layout{60}


\TextSlideTeaser
{Hubschrauberlandeplatz}%1 Titel
{%
Die allgemeinen Vorschriften für
Hubschrauberlandeplätze sind zu beachten.
Der Hubschrauber benötigt einen
geeigneten Landeplatz, i.\,d.\,R. mind.
\EE{25\,$\times$\,25\Meter*}, es dürfen keine
Oberleitungen, lose Gegenstände o.\,ä. den gefahrlosen
Anflug stören, der Untergrund muss tragfähig und nicht
leicht verwehbar sein (Sand!).
Zu Bedenken sind die entstehenden \EE{Luftwirbel}, welche die SanHiSt
oder auch den gesamten Einsatzraum beeinträchtigen
können.
Hier muss ein Kompromiss zwischen den vorhandenen
geeigneten Flächen, der wünschenswerten Nähe und der
erforderlichen Entfernung zur SanHiSt gefunden werden.}%2 Text
{
\begin{itemize}
%     \item Die allgemeinen Vorschriften für
%     Hubschrauberlandeplätze sind zu beachten
    \item Freie ebene Fläche, mind.
    {25\,$\times$\,25\Meter*}
    \item Sichere Umgebung (Stromleitungen, lose Gegenstände, Untergrund, \ldots)
    \item Luftwirbel
\end{itemize}}%3 Slide
{./images/khd/KHD-Symbole-RK/pdf/Transport_heli}%4 Bilddatei
{Kennzeichnung Hubschrauberlandeplatz}%5 Bildtitel
{}%6 Bildbeschreibung
{}%7 Quelle
{}%8 Lizenz
 % Vormals \Layout{60}

% 
% 
% 
% \Layout{22}{Exkurs: Transportmanagement}{%
% \begin{inparaitem}
%     \item Spiralförmige Verteilung
%     \item Zeitliche Staffelung
%     \item Fachliche Zuordnung
%     \item Katastrophe nicht ins Spital verlagern!
% \end{inparaitem}
% }{
% 
% }
% {./images/noauth/khd-transportmanagement.png}
% {}
% {}






\subsection{Sammelstellen}


\TextSlideTeaser
{Sammelstelle Unverletzte}%1 Titel
{%
Hier erfolgt die Betreuung unverletzter Beteiligter. 
Auch Unverletzte bedürfen eines der Witterung angemessenen
Schutzes (Regen, Kälte, \ldots). Mitunter kann die
Beiziehung von Feldküchen o.\,ä. ratsam sein. weiters steht
die Beruhigung, Information und gegebenenfalls die
Beiziehung von geschultem psychologischen Personal
(Kriseninterventionsteam, KIT) im Vordergrund.}%2 Text
{\SymbolText}%3 Slide
{./images/khd/KHD-Symbole-RK/pdf/Unverletzte}%4 Bilddatei
{}%5 Bildtitel
{Kennzeichnung Sammelstelle Unverletzte}%6 Bildbeschreibung
{}%7 Quelle
{}%8 Lizenz
 % Vormals \Layout{60}



\TextSlideTeaser
{Sammelstelle Tote}%1 Titel
{%
 Hier erfolgt die Lagerung der Leichen. Sie
 sollte möglichst diskret eingerichtet werden, abgeschirmt
 vom Behandlungsraum, der Sammelstelle Unverletzte und der
 Außenwelt.  Die Leichen werden abgedeckt, für einen ausreichenden Sichtschutz ist zu sorgen.}%2 Text
{\SymbolText}%3 Slide
{./images/khd/KHD-Symbole-RK/pdf/Tote}%4 Bilddatei
{Kennzeichnung Sammelstelle Tote}%5 Bildtitel
{}%6 Bildbeschreibung
{}%7 Quelle
{}%8 Lizenz
 % Vormals \Layout{60}




\subsection{Organisation}
\label{topic:khd-organisation}


\TextSlideTeaser
{Material- / Meldestelle}%1 Titel
{%
Sie dient zur Registrierung des Personals und Materials. 
Eintreffende Einheiten haben sich bei dieser zu melden. Sie
steht im engen Kontakt zum Leiter SanHiSt und EL.

\vspace{1.5cm}
}%2 Text
{\SymbolText}%3 Slide
{./images/khd/KHD-Symbole-RK/pdf/M}%4 Bilddatei
{Kennzeichnung Material- / Meldestelle}%5 Bildtitel
{}%6 Bildbeschreibung
{}%7 Quelle
{}%8 Lizenz
 % Vormals \Layout{60}


% \Layout{60}{Mobile Leitstelle (MLS)}{%
% \label{topic:mobile-leitstelle}%
% Die \Abk{MLS} beherbergt i.\,d.\,R. die Ein\-satz\-leitung sowie
% wichtige Administrations-, Tele\-kom\-munikations- und
% Funkeinrichtungen. Bei Bedarf werden u.\,U. 
% Handfunkgeräte an die Bereichsleiter ausgegeben. Die MLS ist
% die Brücke zwischen verschiedenen Funkkanälen und zuständig.
% Wenn die MLS die Einsatzleitung beinhaltet, wird sie mit
% einem rotem Dreh- oder Blinklicht gekennzeichnet.
% }{
% {\SymbolText}
% }
% {./images/khd/khd-mls.jpg}
% {}
% {}




\subsection{Führungsstruktur in der SanHiSt}

\TextSlideTeaser
{Einsatzleiter Rettungsdienst}%1 Titel
{%
Ihm obliegt die Organisation aller  rettungsdienstlichen 
Einsatzaufgaben und Rettungsdienst-Einsatzkräfte im gesamten
Schadensraum.
Er ist damit der Vorgesetzte für die eingesetzten Sanitäter und
sonstigen sanitätsdienstlichen Hilfskräfte. Er hat \E{keine}
direkte Weisungsbefugnis gegenüber anderen eingesetzten
Organisationen. Er kann Ärzten \E{keine Weisungen in
medizinischen Fragen} geben.
In der Regel ist der Einsatzleiter RD auch der Leiter der SanHiSt.
\cite{RkRahmenvorschriftGrossunfaelle2007}
}%2 Text
{\SymbolText}%3 Slide
{./images/khd/KHD-Symbole-RK/pdf/EL}%4 Bilddatei
{}%5 Bildtitel
{}%6 Bildbeschreibung
{}%7 Quelle
{}%8 Lizenz
% Vormals \Layout{60}

\TextSlideTeaser
{LNA Leitender Notarzt}%1 Titel
{%
Die Funktion des leitenden Notarztes ist grundsätzlich im Ärztegesetz 
geregelt, 
darüber hinaus gibt es in den Landesgesetzen genauere Regelungen. 
Er ist u.\,a. hinsichtlich der Lagebeurteilung, 
des Sammelns und Sichtens von Verletzten,
der Festlegung von Behandlungsprioritäten,
der medizinische Leitung von Sanitätshilfsstellen,
sowie der Beurteilung des Nachschubbedarfs
und der ärztlichen Beratung der Einsatzleitung 
geschult.
%
Er ist gegenüber den am Einsatz beteiligten Ärzten und 
Sanitätspersonen weisungsbefugt 
und somit der \E{direkte Vorgesetzte aller am Einsatz beteiligten 
Ärzte
\cite{RkRahmenvorschriftGrossunfaelle2007}}.
}%2 Text
{\SymbolText}%3 Slide
{./images/khd/KHD-Symbole-RK/pdf/LNA}%4 Bilddatei
{}%5 Bildtitel
{}%6 Bildbeschreibung
{}%7 Quelle
{}%8 Lizenz
% Vormals \Layout{60}

\TextSlideTeaser
{Leiter SanHist}%1 Titel
{%
Der Lt. SanHiSt organisiert und führt in Absprache mit dem leitenden 
Notarzt die Sanitätshilfsstelle, indem er insbesondere 
die Räume und Stellen festlegt,
das ihm zur Verfügung stehende Personal einteilt,
die Leiter bezeichnet,
das notwendige Material anfordert,
die Verbindung zur Einsatzleitung hält
und eine eventuell notwendige die Ablösung anfordert.
Ihm sind sämtliche Leiter der SanHist-Stellen und Plätze unterstellt, 
ausgenommen den Einsatzleiter RD und den leitenden Notarzt (LNA).
\cite{RkRahmenvorschriftGrossunfaelle2007}
}%2 Text
{
\begin{itemize}
  \item Festlegung Räume und Stellen
  \item Einteilung Personal 
  \item Einteilung Leiter
  \item Materialanforderung
  \item  Verbindung zur Einsatzleitung 
  \item Organisation von Ablösung 
\end{itemize}

}%3 Slide
{./images/khd/KHD-Symbole-RK/pdf/Ltr_SanHist}%4 Bilddatei
{}%5 Bildtitel
{}%6 Bildbeschreibung
{}%7 Quelle
{}%8 Lizenz
% Vormals \Layout{60}

\TextSlideTeaser
{Leiter Material-/Meldestelle}%1 Titel
{%
Er hält engen Kontakt mit Leiter SanHiSt und 
führt Evidenz über eingetroffene Personen und Material. 
\cite{RkRahmenvorschriftGrossunfaelle2007}

\vspace{3\parskip}

}%2 Text
{
\begin{itemize}
  \item Kontakt mit Leiter SanHiSt
  \item Evidenz über eingetroffene Personen und Material
\end{itemize}
}%3 Slide
{./images/khd/KHD-Symbole-RK/pdf/M}%4 Bilddatei
{}%5 Bildtitel
{}%6 Bildbeschreibung
{}%7 Quelle
{}%8 Lizenz
% Vormals \Layout{60}



\TextSlideTeaser
{Leiter Behandlungsraum}%1 Titel
{%
Er betreibt und führt den Behandlungsraum,
insbesonders 
fordert er das notwendige Material an,
legt die Einzelheiten für seinen Raum fest,
lässt die Behandlungsstellen aufbauen und
teilt Personal zu.
Weiters koordiniert er zwischen den Behandlungsstellen.
\cite{RkRahmenvorschriftGrossunfaelle2007}
}%2 Text
{\SymbolText}%3 Slide
{./images/khd/KHD-Symbole-RK/pdf/Behandlung}%4 Bilddatei
{}%5 Bildtitel
{}%6 Bildbeschreibung
{}%7 Quelle
{}%8 Lizenz
% Vormals \Layout{60}




\TextSlideTeaser
{Leiter Transportraum}%1 Titel
{%
Er führt den Transportraum entsprechend der den Patienten zugewiesenen
Transportprioritäten.
Er gliedert seinen Raum, 
indem er die Verladestelle, 
den Hubschrauber-Landeplatz 
und den KFZ- Sammelplatz bezeichnet, 
sowie die Absperrung des Landeplatzes veranlasst.
%
Dabei sorgt er für die ungehinderte Zu- und Wegfahrt im Transportraum
und fordert ggfs.Unterstützung durch die Exekutive an.
%
Er hat einen Überblick über die Aufnahmekapazitäten der Zielspitäler 
und weist den Rettungsfahrzeugen 
und Hubschraubern die Zielspitäler zu.
%
Er führt ein Transportprotokoll 
und hält Evidenz über die zugewiesenen Fahrzeuge.
%
Er setzt bei Notwendigkeit einen Leiter Verladestelle, einen Leiter 
KFZ-Sammelplatz, und
einen Leiter Landeplatz ein.
\cite{RkRahmenvorschriftGrossunfaelle2007}
}%2 Text
{\SymbolText}%3 Slide
{./images/khd/KHD-Symbole-RK/pdf/Transport}%4 Bilddatei
{}%5 Bildtitel
{}%6 Bildbeschreibung
{}%7 Quelle
{}%8 Lizenz
% Vormals \Layout{60}


\TextSlideTeaser
{Leiter Triageraum}%1 Titel
{%
Dieser muss ein Arzt sein.
Er organisiert die Triagestelle(n) 
und teilt allfälligen weiteren Triage-Ärzten die Aufgaben zu.
%
Dabei setzt er auch mobile Triagegruppen ein.
%
Er triagiert soweit möglich selbst auf einer Triagestelle.

Die \E{Triage-Ärzte} sind dem Leiter Triageraum unterstellt 
und sind entweder bei der Bergetriage 
oder in der Triagestelle tätig.
%
Ihre Aufgabe ist die Feststellung der Bergungspriorität
und Festlegen der Triagegruppe (Behandlungspriorität). 
%
Eventuell erteilen sich auch Aufträge für Sanitätshilfe-Maßnahmen.
\cite{RkRahmenvorschriftGrossunfaelle2007}
}%2 Text
{\SymbolText}%3 Slide
{./images/khd/KHD-Symbole-RK/pdf/Ltr_Triage}%4 Bilddatei
{}%5 Bildtitel
{}%6 Bildbeschreibung
{}%7 Quelle
{}%8 Lizenz
% Vormals \Layout{60}


\TextSlideTeaser
{Leiter Information, Informationsstelle}%1 Titel
{%
Die Informationsstelle hat die Aufgabe, 
Behördenvertreter,
Pressevertreter
und sonstige Betroffene (Angehörige, Unverletzte, etc.)
mit notwendigen Informationen zu versorgen und zu betreuen. 
Sich meldende Betroffene werden hier durch Mitarbeiter des 
Kriseninterventionsteams (KIT) empfangen, 
informiert 
und einer entsprechenden Betreuung zugeführt.
\cite{RkRahmenvorschriftGrossunfaelle2007}
}%2 Text
{\SymbolText}%3 Slide
{./images/khd/KHD-Symbole-RK/pdf/Info}%4 Bilddatei
{}%5 Bildtitel
{}%6 Bildbeschreibung
{}%7 Quelle
{}%8 Lizenz
% Vormals \Layout{60}


\section{Patientenleitsystem (PLS)}
\label{topic:pls}

\Text{\TxBeschreibung}{
Das \index{Patientenleitsystem}\T{Patientenleitsystem}
(\index{PLS}\T{PLS}) ermöglicht eine \EE{eindeutige
Kennzeichnung} von Patienten und Zuordnung von
Behandlungsmaßnahmen und Befunden. Der wichtigste
Bestandteil ist die
\index{Patientenleittasche}\T{Patientenleittasche}
(\index{PLT}\T{PLT}).

Sie besteht aus einer orangenen Hülle welche \E{eindeutig
nummeriert} ist. Auf dieser Kunststoffhülle können wichtige
Notizen wie Triagegruppe, etc., vorgenommen werden. In der
Hülle befinden sich das \E{Behandlungsprotokoll} (blau),
\E{Identifikationsprotokoll} (rosa), \E{Zusatzkarten},
\E{Kontaminationsaufkleber} und ein Bogen mit gleichlautend
\E{nummerierten Klebeetiketten}. Am unteren Ende befinden
sich \E{zwei Abschnitte} welche die Rückverfolgbarkeit des
Transports sicherstellen. Die PLT ist mit einer Gummischnur
versehen damit sie dem Patienten umgehängt werden kann.

\EE{\E{Spätestens} an der Triagestelle} bekommt jeder
Betroffene \enquote{seine} \Abk{PLT}. Auch Unverletzte und
Verstorbene werden mittels \Abk{PLS} registriert.
}{}{}{}{}

\begin{Synopsis}
  \SynopsisItem{Patientenleittasche: 
  Die wichtigste Funktion ist die 
  \EE{eindeutige Registrierung} der Patienten.}
\end{Synopsis}

\Text{Zusatzkarten}{%
Bei einer eventuellen Bergetriage bekommt jeder Patient eine
Patientenleittasche (\Abk{PLT}). Patienten, welche
\EE{bevorrangt gerettet} werden sollen, werden zusätzlich
mit der gelben {\fcolorbox{Black}{Yellow}{\sffamily\bfseries
DRINGEND}}\,-Karte gekennzeichnet. Die Karte wird auf der
Triagestelle wieder entfernt. Tote erhalten zusätzlich zur
\Abk{PLT} eine schwarze
{\fcolorbox{Black}{White}{\sffamily\bfseries
VERSTORBEN}}\,-Karte.
}{}{}{}{}

\Text{Behandlungsprotokoll / Identifikationsprotokoll}{%
Die Protokolle werden im Behandlungsraum ausgefüllt wenn
Zeit dafür ist. Das Ausfüllen der Protokolle darf den
Abtransport unter keinen Umständen verzögern. Die Protokolle
können auch leer bleiben.
}{}{}{}{}

\Text{Kontaminationsaufkleber}{%
\index{Kontaminationsaufkleber}%
Wenn der Patient kontaminiert ist, wird ein
Kontaminationsaufkleber in Form eines gelben Dreiecks mit
schwarzem Rand auf die \Abk{PLT} geklebt. Dies soll alle
nachfolgenden Personen, welche mit dem Patienten in Kontakt
treten, warnen. Die Aufkleber sind auf \EE{beiden Seiten}
der \Abk{PLT} anzubringen!
}{}{}{}{}

\Text{PLT-Nummer}{%
Die \Abk{PLT} und damit der Patient ist durch eine
\EE{eindeutige, einmalige Nummer} \EE{identifiziert}.
Mittels \EE{Selbstklebeetiketten} können die Nummern der
Patienten in verschiedenen Listen schnell erfasst werden.
Das ist in vielen Bereichen besonders wichtig (Triage,
Behandlung, Transport). Außerdem können alle Effekte des
Patienten damit gekennzeichnet werden.

Die Patientenleittasche verbleibt solange am Patienten,
\E{bis das Aufnahmeverfahren im Krankenhaus abgeschlossen
ist}, d.h. bis der Patient im Spital administrativ erfasst
ist. Anschließend wird die \Abk{PLT} der Krankengeschichte
des Patienten beigefügt.
}{}{}{}{}

\begin{Synopsis}
  \SynopsisItem{Die Nummerierung der \Abk{PLT} ermöglicht die
  eindeutige Zuordnung und Erfassung von
  Patienten. Das ist eine der Hauptaufgaben der \Abk{PLT}!}
\end{Synopsis}

\noindent\begin{minipage}{\linewidth}
\noindent\includegraphics[width=\linewidth]{./images/khd/plt-inhalt-beschriftet.jpg}

\begin{multicols}{2}

\begin{description}
    \item[a] Zusatzkarte {\fcolorbox{Black}{White}{\sffamily\bfseries
VERSTORBEN}}
    \item[b] Zusatzkarte 
{\fcolorbox{Black}{Yellow}{\sffamily\bfseries DRINGEND}}
    \item[c] Klebeetiketten mit PLT-Nummer
    \item[d] Aufkleber für kontaminierte
Gegenstände 
    \item[e] Behandlungsprotokoll (blau)
    \item[f] Identifikationsprotokoll (rosa)
\end{description}

\end{multicols}

\end{minipage}




\begin{PictureSeriesWide}{}{Übersicht: \SlideHeading{Felder der Patientenleittasche}}
\PictureSeriesRowWide{4}
{./images/khd/plt-vorderseite-beschriftet-50pc.jpg}
{Das Feld \EE{Triage} \textcircled{a} dient der  
Kennzeichnung der Triagegruppe ({\sffamily I}, {\sffamily II},
{\sffamily III} oder {\sffamily IV}). Der \EE{zweite Triageblock} \textcircled{b}
ermöglicht die  Umtriage oder Spitalstriage. Das Feld
\EE{Transport} \textcircled{c} kennzeichnet die \E{Transportpriorität} 
({\sffamily A}: hoch; {\sffamily B}: niedrig), diese wird später
entschieden und ist unabhängig von der Triagegruppe.
}
{./images/khd/plt-vorderseite-umtriagierung15-50pc.jpg}
{PLT: Umtriagierung innerhalb der ersten
Viertelstunde: Mittels Pfeil vom vorigen zum nun zutreffenden Triagefeld in der gleichen Zeile}
{./images/khd/plt-vorderseite-umtriagierung60-50pc.jpg}
{PLT: Umtriagierung nach längerer Zeit: Neues Kreuz beim entsprechenden Triagefeld in der unteren Zeile}
{./images/khd/plt-rueckseite-angekreuzt-50pc.jpg}
{Der Triagearzt trägt auf der Rückseite die zu setzenden
\EE{Ma{\ss}nahmen} ein. Der Durchführende (Arzt, NFS, RS) bestätigt
dann die Durchführung der
Maßnahmen im rechten Feld durch eine Paraphe.
}
\end{PictureSeriesWide}

\TextSlide
{Abrisse}
{
Die Patientenleittasche verfügt über 2 Abrisse:

\begin{enumerate}
  \item \EE{Abriss für die SanHiSt} -- Leiter Transport: 
  Dieser verbleibt in der SanHiSt beim Leiter Transport, sodass nachvollzogen werden 
  kann, welche Patienten bereits abtransportiert wurden. Auf dem Abschnitt 
  wird die Verladezeit, die KFZ-Nummer, das Zielspital und die Zielabteilung angegeben.
  \item \EE{Abriss für das Zielspital}: Dieser wird im Zielspital 
  abgegeben und ermöglicht dort eine administrative Erfassung des Patienten. 
  Der Abschnitt wird mit der Ankunftszeit versehen.
\end{enumerate}
}
{
\begin{itemize}
  \item {Abriss für die SanHiSt -- Leiter Transport}
  \item {Abriss für das Zielspital}
\end{itemize}}
{}
{}
{}



\section{Spezielles Material}
\label{topic:katastrophe-material}

Zur Bewältigung von Großschadenslagen werden spezielle
Materialien und Einheiten von diversen Organisationen
vorgehalten. Sie sind hinsichtlich ihrer genauen Ausstattung
nicht überregional normiert, sondern den jeweiligen
Erfordernissen angepasst. Die wichtigsten sind:

\begin{description}
  \item[\II{BaGUS}] \I{Basis-Großunfallset}: Materialien zur
  sanitätsdienstlichen Versorgung, z.\,B.
  Schienungsmaterial, Decken, Tragen, Bergetücher {\ldots} Der genaue
  Inhalt ist nicht genormt. 
  \item[\II{BeGUS}] \I{Beleuchtungs-Großunfallset}:
  Beleuchtungskörper und Zubehör zur Ausleuchtung von
  Behandlungs- oder Schadensplätzen. Der genaue
  Inhalt ist nicht genormt.
  \item[\II{MeGUS}] \I{Medizinisches Großunfallset}: Es
  enthält Materialien zur erweiterten medizinischen Versorgung von
  Patieten, z.\,B. Medikamente, Infusionen. Der genaue
  Inhalt ist nicht genormt.
  \item[\II{MobSAN}] Mobile Sanitätstruppe, bestehend aus
  einer bestimmten Anzahl von RTWs, KTWs und anderen Fahrzeugen
  \item[\II{MobSAN-A}] Wie MobSAN, jedoch mit Arzt und NAW
  oder NEF
  \item[Versorgungseinheiten] wie z.\,B. Feldküchen sind bei
  längeren Szenarien wichtig, um die längerfristige
  Einsatzbereitschaft sicherzustellen. 
  \item[Zelte] gibt es in den verschiedensten Ausführungen,
  aufgrund des relativ aufwändigen Aufbaus eignen sie sich
  eher bei länger dauernden Schadensfällen
\end{description}

\subsection{Sonderfall: Großambulanzdienste}

Bei der sanitätsdienstlichen Betreuung von
Großveranstaltungen finden sich zuweilen ähnliche Probleme
wie im Großschadenfall. Daher bedient man sich oft einiger
der dort angewendeten Techniken, wobei jedoch der
Schwerpunkt bei der Individualmedizin liegt, z.\,B.:

\begin{itemize}
  \item Gliederung der SanHiSt (Wagenhaltplätze,
  Behandlungsräume, \ldots)
  \item Führungsstruktur (Material- und Meldestelle, div.
  Leiter \ldots)
  \item Material (mobile Leistellen, MeGUS- oder
  BaGUS-Einheiten, \ldots)
\end{itemize}
 
Im Gegensatz zum
Großschadensfall werden Großereignisse vorab geplant,
sodass für eine entsprechende Infrastruktur und
ausreichende Ressourcen gesorgt werden kann.
Dementsprechend obliegt es auch den mit der Organisation
 Betrauten festzulegen, welche Verfahrensweisen im
jeweiligen Fall zur Anwendung kommen.

\begin{Synopsis}
  \SynopsisItem{Großveranstaltungen sind im Gegensatz zur
  Großschadensfällen \EE{planbar}!}
\end{Synopsis}
 

\ChapterEnd
